\chapter{Comparatives and Superlatives}
\label{chap:comparison}

Comparison describes a difference or places something at the top of a group.
Learn the German names alongside the English ones: \textit{Positiv}
(basic form), \textit{Komparativ} (comparative), and \textit{Superlativ}
(superlative).

\section{The three degrees}

\begin{center}
\small
\rowcolors{2}{referenceBlueBg}{white}
\begin{tabularx}{\textwidth}{@{}lllX@{}}
\toprule
Degree & Form & Meaning & Example \\
\midrule
Positiv & schnell & fast & \textit{Der Zug ist schnell.} \\
Komparativ & schneller & faster & \textit{Der Zug ist schneller als der Bus.} \\
Superlativ & am schnellsten & fastest & \textit{Von diesen Zügen ist dieser am schnellsten.} \\
\bottomrule
\end{tabularx}
\end{center}

The comparison group may be stated or understood. A comparative can stand
alone when the context supplies the comparison: \textit{Dieses Zimmer ist
größer} (this room is bigger).

\section{Building the forms}

The usual comparative ending is \textbf{-er}; the usual superlative stem
ends in \textbf{-st} or \textbf{-est}. With \textit{am}, the complete form
ends in \textbf{-sten} or \textbf{-esten}. German uses these endings even
with long adjectives: \textit{interessant -- interessanter -- am interessantesten}.

\begin{center}
\small
\rowcolors{2}{referenceGreenBg}{white}
\begin{tabularx}{\textwidth}{@{}l l l X@{}}
\toprule
Pattern & Positiv & Komparativ & Superlativ with \textit{am} \\
\midrule
Regular & schön & schöner & am schönsten \\
Umlaut & jung & jünger & am jüngsten \\
Extra \textit{e} & laut & lauter & am lautesten \\
Unstressed \textit{e} drops & dunkel & dunkler & am dunkelsten \\
Shorter comparative stem & teuer & teurer & am teuersten \\
Special superlative spelling & groß & größer & am größten \\
\bottomrule
\end{tabularx}
\end{center}

Some short adjectives change \textit{a, o, u} to \textit{ä, ö, ü} in both
comparison forms: \textit{alt -- älter}, \textit{groß -- größer},
\textit{kurz -- kürzer}. Learn this with each word: an eligible vowel does
not guarantee an umlaut, as \textit{klar -- klarer} shows.
The extra \textit{e} in a superlative can help pronunciation:
\textit{heiß -- am heißesten}, \textit{alt -- am ältesten}.
The reference tables below give complete forms.

\begin{tcolorbox}[colback=referenceYellowBg,colframe=genderPlural,title={Forms to learn directly},fonttitle=\bfseries]
\textit{gut -- besser -- am besten}\qquad
\textit{hoch -- höher -- am höchsten}\par
\textit{nah -- näher -- am nächsten}\qquad
\textit{groß -- größer -- am größten}
\end{tcolorbox}

\section{Comparisons in sentences}

Use \textbf{als} after a comparative in standard German. Equality uses
\textbf{so/genauso \ldots{} wie} with the basic form.

\begin{center}
\small
\rowcolors{2}{referenceBlueBg}{white}
\begin{tabularx}{\textwidth}{@{}lXX@{}}
\toprule
Meaning & Pattern & Example \\
\midrule
More of a quality & Komparativ + als
  & \textit{Die Bahn ist schneller als der Bus.} \\
Equal degree & so / genauso + Positiv + wie
  & \textit{Mein Zimmer ist genauso groß wie dein Zimmer.} \\
Lower degree & weniger + Positiv + als
  & \textit{Dieser Weg ist weniger steil als der andere.} \\
Below the same degree & nicht so + Positiv + wie
  & \textit{Heute ist es nicht so kalt wie gestern.} \\
Increasing degree & immer + Komparativ
  & \textit{Die Tage werden immer länger.} \\
Size of a difference & viel / etwas + Komparativ
  & \textit{Das ist viel billiger. Dieser Weg ist etwas kürzer.} \\
Highest degree in a group & am + Superlativ
  & \textit{Von den drei Wegen ist dieser am kürzesten.} \\
\bottomrule
\end{tabularx}
\end{center}

\textit{Billiger} already means ``cheaper''; \textit{viel billiger} means
``much cheaper''. A regular comparative does not need an additional
\textit{mehr}.

\subsection{Two linked changes: \textit{je \ldots{} desto/umso}}

\textit{Je mehr ich übe, desto sicherer spreche ich.}
\quad--- The more I practise, the more confidently I speak.

The \textit{je}-clause puts its finite verb at the end. In the main clause,
\textit{desto/umso + comparative} forms the first unit; the finite verb
comes immediately after it. \textit{Umso} can replace \textit{desto}:
\textit{Je länger wir warten, umso später kommen wir an.}

\section{Using comparison forms before a noun}

After \textit{sein} or when describing an action, the comparative has no
case ending: \textit{Das Auto ist schneller. Sie fährt schneller.}
Before a noun, add the adjective ending required by the article, case,
gender and number. The same ending tables apply at every degree; see
Section~\ref{sec:adjective-endings} on page~\pageref{sec:adjective-endings}.

\begin{center}
\small
\rowcolors{2}{referenceGreenBg}{white}
\begin{tabularx}{\textwidth}{@{}lXXX@{}}
\toprule
Use & Positiv & Komparativ & Superlativ \\
\midrule
After \textit{sein} & \textit{Das Auto ist schnell.}
  & \textit{Das Auto ist schneller.}
  & \textit{Das Auto ist am schnellsten.} \\
With a definite article & \textit{das schnelle Auto}
  & \textit{das schnellere Auto}
  & \textit{das schnellste Auto} \\
\bottomrule
\end{tabularx}
\end{center}

In \textit{ein schnelleres Auto}, the comparative \textit{schneller}
receives the neuter ending \textit{-es}. In \textit{mit dem schnellsten Auto},
the superlative stem \textit{schnellst-} receives \textit{-en}.
Choose the noun's case from its role in the sentence, not from the comparison.

\begin{tcolorbox}[colback=referenceYellowBg,colframe=genderPlural,title={Two useful superlative patterns},fonttitle=\bfseries]
\textit{Dieser Berg ist am höchsten.}\quad--- This mountain is the highest.\par
\textit{Das ist der höchste Berg der Region.}\quad--- That is the highest mountain in the region.\par
Use \textit{am höchsten} without a following noun, or an inflected form
such as \textit{der höchste Berg} before a noun. The endings change with the
case: \textit{auf dem höchsten Berg}.
\end{tcolorbox}

\section{Common adjective forms}

These tables show one common form for each degree. The \textit{am}-form
is a convenient way to learn the superlative; its use before a noun follows
the ending rules above.

\noindent
\begin{minipage}[t]{0.49\textwidth}
\vspace{0pt}
\footnotesize
\setlength{\tabcolsep}{2pt}
\renewcommand{\arraystretch}{1.12}
\rowcolors{2}{referenceBlueBg}{white}
\begin{tabularx}{\linewidth}{@{}lllX@{}}
\toprule
English & Positiv & Komparativ & Superlativ \\
\midrule
old & alt & älter & am ältesten \\
pleasant & angenehm & angenehmer & am angenehmsten \\
tense / strained & angespannt & angespannter & am angespanntesten \\
poor & arm & ärmer & am ärmsten \\
cheap & billig & billiger & am billigsten \\
wide & breit & breiter & am breitesten \\
thick / fat & dick & dicker & am dicksten \\
stupid & dumm & dümmer & am dümmsten \\
dark & dunkel & dunkler & am dunkelsten \\
thin & dünn & dünner & am dünnsten \\
easy / simple & einfach & einfacher & am einfachsten \\
narrow / tight & eng & enger & am engsten \\
relaxed & entspannt & entspannter & am entspanntesten \\
hard-working & fleißig & fleißiger & am fleißigsten \\
friendly & freundlich & freundlicher & am freundlichsten \\
fresh & frisch & frischer & am frischesten \\
early & früh & früher & am frühesten \\
dangerous & gefährlich & gefährlicher & am gefährlichsten \\
healthy & gesund & gesünder & am gesündesten \\
happy / lucky & glücklich & glücklicher & am glücklichsten \\
big / tall & groß & größer & am größten \\
good & gut & besser & am besten \\
hard & hart & härter & am härtesten \\
frequent & häufig & häufiger & am häufigsten \\
hot & heiß & heißer & am heißesten \\
bright & hell & heller & am hellsten \\
high & hoch & höher & am höchsten \\
interesting & interessant & interessanter & am interessantesten \\
young & jung & jünger & am jüngsten \\
cold & kalt & kälter & am kältesten \\
\bottomrule
\end{tabularx}
\end{minipage}%
\hfill
\begin{minipage}[t]{0.49\textwidth}
\vspace{0pt}
\footnotesize
\setlength{\tabcolsep}{2pt}
\renewcommand{\arraystretch}{1.12}
\rowcolors{2}{referenceBlueBg}{white}
\begin{tabularx}{\linewidth}{@{}lllX@{}}
\toprule
English & Positiv & Komparativ & Superlativ \\
\midrule
clear & klar & klarer & am klarsten \\
small & klein & kleiner & am kleinsten \\
clever & klug & klüger & am klügsten \\
ill / sick & krank & kränker & am kränksten \\
short & kurz & kürzer & am kürzesten \\
long & lang & länger & am längsten \\
slow & langsam & langsamer & am langsamsten \\
loud & laut & lauter & am lautesten \\
light / easy & leicht & leichter & am leichtesten \\
tired & müde & müder & am müdesten \\
near & nah & näher & am nächsten \\
wet & nass & nässer & am nässesten \\
new & neu & neuer & am neuesten \\
rich & reich & reicher & am reichsten \\
quiet / calm & ruhig & ruhiger & am ruhigsten \\
sour & sauer & saurer & am sauersten \\
bad & schlecht & schlechter & am schlechtesten \\
beautiful & schön & schöner & am schönsten \\
fast & schnell & schneller & am schnellsten \\
weak & schwach & schwächer & am schwächsten \\
heavy / difficult & schwer & schwerer & am schwersten \\
safe / certain & sicher & sicherer & am sichersten \\
exciting & spannend & spannender & am spannendsten \\
strong & stark & stärker & am stärksten \\
expensive & teuer & teurer & am teuersten \\
dry & trocken & trockener & am trockensten \\
warm & warm & wärmer & am wärmsten \\
soft & weich & weicher & am weichsten \\
far / wide & weit & weiter & am weitesten \\
important & wichtig & wichtiger & am wichtigsten \\
\bottomrule
\end{tabularx}
\end{minipage}
\par\smallskip

The forms shown for \textit{gesund} and \textit{nass} use an umlaut;
\textit{gesunder -- am gesundesten} and \textit{nasser -- am nassesten}
are also correct. In \textit{dunkler}, \textit{teurer} and \textit{saurer},
notice the shorter comparative stem.

\subsection{From verbs to adjectives: tension and relaxation}

\textit{Der Film ist spannend. Vor der Prüfung bin ich angespannt.
Nach dem Urlaub bin ich entspannter.}
These adjectives share a verb family but describe different qualities.

\clearpage
\section{Quantities, preferences and limits}

Quantity words and a few adverbs also have comparison forms. Learn these
separately from the ordinary adjective endings.

\begin{center}
\small
\rowcolors{2}{referenceYellowBg}{white}
\begin{tabularx}{\textwidth}{@{}llllX@{}}
\toprule
Meaning & Basic form & Comparative & Superlative & Example \\
\midrule
much / many & viel & mehr & am meisten & \textit{Heute habe ich mehr Zeit.} \\
little / few & wenig & weniger & am wenigsten & \textit{Dieses Ticket kostet am wenigsten.} \\
gladly & gern(e) & lieber & am liebsten & \textit{Ich fahre am liebsten mit dem Zug.} \\
\bottomrule
\end{tabularx}
\end{center}

With nouns, compare \textit{mehr Zeit} and \textit{die meiste Zeit},
or \textit{weniger Arbeit} and \textit{die wenigste Arbeit}.
For preferences, \textit{Ich lese lieber} means ``I prefer reading'';
\textit{Ich lese am liebsten} means ``I like reading best''.

Not every adjective expresses a quality with a measurable degree.
Words such as \textit{schriftlich} or \textit{monatlich} normally classify
something, so ordinary literal uses do not call for comparison forms.
An adjective's meaning and context determine whether comparison is useful;
figurative and playful uses can behave differently.
